\documentclass[11pt]{article}

\usepackage[margin=1in]{geometry}
\usepackage{amsmath,amssymb}
\usepackage{hyperref}
\usepackage{enumitem}

\hypersetup{
  colorlinks=true,
  linkcolor=blue,
  urlcolor=blue,
  citecolor=blue
}
\emergencystretch=2em

\newcommand{\omegaone}{\omega_1}
\newcommand{\Natural}{\mathbb{N}}

\title{A Literature-Implied Countable-Compacta Obstruction for Gorak's Net Banach-Stone Problem}
\author{Math discovery packet}
\date{June 30, 2026}

\begin{document}
\maketitle

\section{Status}

This is a lightweight \emph{literature-implied answer}, not an original
solution.  The source paper is R. Gorak, \emph{Coarse version of the
Banach-Stone theorem}, arXiv:1005.0937.
In the final remarks, Problem \texttt{problem Banach-Stone ogolny} asks
whether, for compact spaces \(X,Y\), the inequality
\[
  d_N(C(X),C(Y))<2
\]
forces \(X\) and \(Y\) to be homeomorphic.  The next problem asks for the
answer when \(X,Y\) are countable compacta, and also asks what happens when
\(X\) is the usual convergent sequence.

The supporting paper is A. Prochazka and L. Sanchez-Gonzalez, \emph{Low
distortion embeddings between \(C(K)\) spaces}, arXiv:1408.0211.

\section{Identification}

The supporting paper recalls the Mazurkiewicz--Sierpinski classification:
every countable Hausdorff compact space is homeomorphic to an ordinal interval
of the form
\[
  [0,\omega^\alpha\cdot n],
  \qquad \alpha<\omegaone,\quad 1\leq n<\omega.
\]

Corollary \texttt{c:AnswerGorak} in arXiv:1408.0211 states that if
\(\gamma\neq \alpha<\omegaone\) and \(m,n\in\Natural\), then
\[
  d_N\!\left(C([0,\omega^\gamma\cdot n]),
             C([0,\omega^\alpha\cdot m])\right)\geq 2.
\]
Immediately after the corollary, the authors state that this partially answers
Problem 2 in Gorak's paper.

Consequently, for countable compacta, net distance strictly below \(2\)
forces equality of the ordinal exponent in the
Mazurkiewicz--Sierpinski normal form.  Equivalently, it rules out all
countable compacta of different Cantor--Bendixson height in the threshold
\(<2\) problem.

\section{Scope}

This does not settle Gorak's countable-compacta problem in full.  The
supporting paper itself notes that the same-exponent finite multiplicity
question remains outside the stated corollary: the spaces
\[
  [0,\omega^\alpha\cdot m]
  \quad\text{and}\quad
  [0,\omega^\alpha\cdot n]
\]
with \(m\neq n\) are not separated by Corollary \texttt{c:AnswerGorak}.

In particular, for the source paper's test case where \(X\) is a convergent
sequence, this packet records only the height obstruction.  It does not prove
that a countable compactum \(Y\) with \(d_N(C(X),C(Y))<2\) must be homeomorphic
to \(X\).

\section{Search Evidence}

The run's lightweight indexes had no existing direct packet for arXiv:1005.0937
or this countable-compacta height obstruction.  The identification comes from
the local source of arXiv:1408.0211, whose introduction and Corollary
\texttt{c:AnswerGorak} explicitly connect the result to Gorak's Problem 2.

\section{References}

\begin{itemize}[leftmargin=*]
\item R. Gorak, \emph{Coarse version of the Banach-Stone theorem},
arXiv:1005.0937.
\item A. Prochazka and L. Sanchez-Gonzalez, \emph{Low distortion embeddings
between \(C(K)\) spaces}, arXiv:1408.0211.
\item S. Mazurkiewicz and W. Sierpinski, classical classification of countable
compact Hausdorff spaces, as cited in arXiv:1408.0211.
\end{itemize}

\end{document}
